\chapter{Case study 2 - quantum physics}

The second case study of a particular model of physics is indeed, quantum physics. Despite such success in modelling and analysing the world, classical physics hits deadlock where behaviours and observations are so little, effects on macroscopic scale are so large, and in some case the system is too small to isolate such that results and predictions come in fruitless and of particular failure. Of such, solution to such problem comes in the form of creating a new theory, connected still in the sense of certain special cases or edge cases, and of removing factors, with \textit{new postulate} and structure accompanied, and of models that would be fairly difficult to incorporate, aside from of modern classical physics. Such is similar to how relativity is done, just with a generalization of missing content. 

The mathematical technique for modelling quantum physics is also fairly different from classical physics, mainly in the way that the quantum mechanical system is treated and how is that formed. Specifically, one will expect to master of linear algebra, mathematical analysis, topology - especially Hilbert space and rigged Hilbert space topology, differential geometry (kinda), manifold analysis, operator theory, tensor theory, probability theory, and group theory, especially Galois group. Such is not even enough for certain niche covering of quantum physics, but that is for the edge case interpretation they require to express such system of interest (not like I want to discuss of field theory right from the start). For said motion, we can concur on an observation that quantum system behaves toward a much, very much different axioms set, and said thing is what we need to cover and analyse of our own interest. 

\section{Mathematical preliminaries}

There are certainly some specification of knowledge when it comes to understanding quantum mechanics. For now, those mathematical concepts might look familiar to somewhere else, but they contain little to no physics of such. Toward that end, let us take an inquiry on whatever the heck is required to understand this eldritch horror.
\section{Linear and inner vector spaces}

By itself, a \textit{vector space} is a linear space. That is why it is even the main structure of the linear algebra course. Of course, again, there exist the nonlinear manifold, the \textit{affine space} and so on, but they are special case of dropping certain linearity criteria. We would then review such concept here for reference purpose. 

\begin{definition}[(Linear) vector space]
    A vector space $\mathbb{V}$ is a set of objects called \textit{vectors}, denoted $v,\vec{v},\mathbf{v}$ or $\ket{v}$, over a field $\mathbf{F}$, equipped with two operators $(+,\cdot)$ for addition and scalar multiplication, such that the following holds: 
    \begin{enumerate}[topsep=1pt,itemsep=1pt,leftmargin=1cm]
        \item Closure: if $V,W\in \mathbb{V}$, then $V+W\in\mathbb{V}$. 
        \item Closure of scalar multiplication: if $V\in \mathbb{V},\lambda\in \mathbb{F}$, then $\lambda V\in \mathbb{V}$. 
        \item Addition commutativity: if $V,W\in \mathbb{V}$, then $V+W=W+V$. 
        \item Addition associativity: if $V,W\in \mathbb{V}$, then $V+W+Z=(V+W)+Z$.
        \item Scalar multiplication associativity: if $a\in \mathbb{F}$ and $V\in\mathbb{V}$, then $a(bV)=(ab)V$. 
        \item Scalar multiplication distributivity: if $a,b\in\mathbb{F}$, and $V,W\in\mathbf{V}$, then $(a+b)V=aV+bV$ and $a(V+W)=aV+aW$. 
        \item Existence of nullity: There exists a \textit{null vector}, denoted $0,\vec{0},\mathbf{0}$ or $\ket{0}$ such that $V+0=V$ and $0V =0$, for $V\in\mathbb{V}$. 
        \item Existence of negatives: There exists for every $v\in\mathbb{V}$ a vector, written as $-v$, such that $v+(-v)=0$. 
        \item Unitary: There exists the identity vector, denoted $1,\mathbf{1},\ket{1}$, such that $1v=v$ for all $v\in\mathbb{V}$. 
    \end{enumerate}
\end{definition}
Here, we make use of a new notation, which stems from the `official' Dirac braket notation, $\ket{v}$ to denote such vector with name. Think of it as an abstract object in such space, and we now only write it like that so it would be familiar later on. for any mathematical structure that is embedded and structured to take form of a vector space, then each individual vector represents certain subject of said structure, for example, the \textit{state} of the structure in particular configuration. If the field itself is replaced of the real number $\mathbb{R}$, then we say it as a real vector space. Otherwise, if the field is the complex field, $\mathbb{C}$, then it is said to be a complex vector space. This topic should be treated on its own section rather than leaving it only here. 

Such encoding allows for two types of general interest vector space. One is the \textit{discrete basis coordinate space}, for general vectors like $v=(x_1,x_2,x_3,\dots)$, or it can be the space of certain \textit{functions} of some type, for example the space of all differentiable functions. In such case, the field still holds as scalar, but now, we have the entire function as a particular basis or coordinate. Usually, in a very mathy term, it is called a \textbf{function space}, and it is a topological space whose points are now functions instead, more like a wrapper above the typical structure. 

Next, our definition of vector space consequently leave us with a construction of a \textit{basis}. Basically, this mean that they possess bases such that for a vector space $E$, there exist some (one or more) set of vectors $\{e_{1},e_{1},\dots,e_{n}\}$ of the same dimension, such that every vector $\ket{v}\in E$ can be represented by 
\begin{equation}
    \ket{v} = \lambda_{1}e_1 + \lambda_{2}e_{2} + \dots + \lambda_{n} e_{n}
\end{equation}
for all $(\lambda_{1},\dots,\lambda_{n}\in \mathbb{R})$. This is defined as the \textit{linear combination} of such basis upon to express the vector $\ket{v}$. We note that this works for, in general, the basis of real space. The same happens with complex space, but is much, much more difficult to analyze. This definition of a basis, on such is valid for \textit{finite-dimensional space}. We will see how we can extend such notion to the case of infinitely dimensional space later on. For a set of vectors, now we might as well denote here of this physics section as $\ket{\cdot}$, so $\ket{1},\ket{2},\dots,\ket{n}$, is said to be used to construct a \textit{linear combination} on vector space $E$ if , denoting $(a_{i})_{i\in I}$ for index set $I$ the \textit{$I$-indexed family} of element of a set $A$ to be the function $a:I\to A$, the \textit{subfamily} of $(u_{i})_{i\in I}$ being the family $(u_{j})_{j\in J}$ for any $J\subset I$, satisfying the following: 
\begin{definition}[Linear combination]
    Let $E$ be a vector space. A vector $v\in E$ is a \textit{linear combination of a family} $(u_{i})_{i\in I}$ of elements of $E$ iff there is a family $(\lambda_{i})_{i\in I}$ of scalars in $\mathbb{R}$ such that 
    \begin{equation}
        v = \sum_{i\in I} \lambda_{i} u_{i}
    \end{equation}
    When $I=\varnothing$, we stipulate that $v=0$. 
\end{definition}
Then, from such definition, we can then define \textit{linear (in)dependence} of such vector set. 
\begin{definition}[Linearly (in)dependent]
    We say that a family $(u_{i})_{i\in I}$ is \textit{linearly independent} iff for every family $(\lambda_{i})_{i\in I}$ of scalar in $\mathbb{R}$, 
    \begin{equation}
        \sum_{i\in I} \lambda_{i}u_{i} = 0 \Rightarrow \lambda_{i} \ne 0 \: \forall i\in I 
    \end{equation}
    That is, the linear set is unique such that only the null vector yields such correspondence. Equivalently, a family $(u_{i})_{i\in I}$ is said to be \textit{linearly dependent} iff there is some family $(\lambda_{i})_{i\in I}$ of scalars in $\mathbb{R}$ such that 
    \begin{equation}
        \sum_{i\in I} \lambda_{i} u_{i} = 0 \Rightarrow \lambda_{j} \neq 0 , j \in I
    \end{equation}
    We agree also that for $I = \varnothing$, then the family $\varnothing$ is linearly independent. 
\end{definition}

The above simply state that it is linear independent iff either $I\neq \varnothing$, or $I$ consists of a single element $i$ and $u_i \neq 0$, or $|I|\leq 2$ and no vector $u_{j}$ in the family can be expressed as a linear combination of other vectors in the family. The reverse happen for linear dependent. Thus, we can see that, for something to be a basis, it needs to be linearly independent. However, we also look at the criteria that any vector $v\in E$ should be able to be expressed by such linear combination. Then, the set of linear independent vectors must also \textit{span} the vector space $E$, or in simpler term, of the same dimension as vector space $E$ indict upon. We then also want to express now the late definition of what constitute the \textit{dimension} of such space. 

\begin{definition}[Dimension]
    A vector space has \textit{dimension} $n$ if it can accommodate a maximum of $n$ linearly independent vectors. For a vector space $V$, this is denoted by $V^{n}(\mathbb{F})$ follow by their field, which is not necessarily real or complex. 
\end{definition}

Moving on with vector space, without the usual vector subspaces and stuff (which has been covered in its respective chapter on linear algebra), we can add additional structure to the space itself. One of such is the \textit{inner product}, or the scalar product for a linear vector space. The inner product is oddly similar to the dot product in the usual expression of vector algebra, such is $\langle A , B \rangle = |A||B|\cos{(\theta)}$. This expression, in coordinate has the form of $\langle A , B \rangle = A_{x}B_{x}+A_{y}B_{y}$ and else, which can be generalized to $n$-space. The formula then indeed specify that dot product as a functional $\langle \cdot, \cdot \rangle : V\times V \to \mathbb{F}$ of the embedding it receives. To this end, it satisfies the symmetry property, i.e. $\langle A , B \rangle = \langle B,A \rangle $, positive semidefiniteness\footnote{The term here might be confusing, but this might help. For a specific function, then, for now is the inner product, being positive definite means that $\langle A , B \rangle >0$ whenever $A\neq 0$. Positive \textit{semidefinite} means that now $\langle A ,A \rangle \geq 0$ instead for all $A$, and could be zero even if $A\neq 0$. In matrix notation, this means you indict that $x^{\top}Ax\geq 0$ for all $x$, and that $A$ symmetric matrix is what called an $A$ positive semidefinite matrix. The only reason that it is there, is because the inner product will ultimately encode the more general notion of dot product, which makes it introducing new geometry on top. Simple case, $A=I$ of identity matrix on $n$-space.}, which is $\langle A , A \rangle \geq 0$, being $0$ iff $A=0$ for $A\in E$, and linearity property 
\begin{equation*}
    \langle A , (\lambda B + \psi C) \rangle = \lambda \langle A , B \rangle + \psi \langle A , C \rangle  
\end{equation*} 
The generalization of such concept is then defined to be the \textit{inner product}. We denote such with a few notations, namely $(\phi,\psi)$, the usual dot product will suffice, $\langle a,b \rangle $ or in a more quantum-ish notation, $\bra{a}\ket{b}$. The following are rules that it must obey.
\begin{align}
    \bra{a}\ket{b} &= \bra{b}\ket{a} \quad \text{(Symmetry)}\\ 
    \bra{a}\ket{a} & \geq 0 , \forall \ket{a} \neq \ket{0}\quad \text{(Positive semidefiniteness)}\\
    \bra{a} (\lambda\ket{W}+\psi\ket{c}) &\equiv \bra{a}\ket{\lambda b + \psi c} = \lambda\bra{a}\ket{b} + \psi\bra{a}\ket{c} \quad \text{(Linearity of the second argument)}
\end{align}
A vector space with an inner product is then called the \textbf{inner product space}\index{inner product space}. We note that the first argument of the inner product is often expressed to be \textit{antilinear}, that is $\bra{u}\ket{v}=(\bra{v}\ket{u})^{*}$ of the conjugate symmetry. In case of real number, since every real number equals its complex conjugate, then $\bra{u}\ket{v}=\bra{v}\ket{u}$ for $u,v\in E$. Any operation that satisfies this can be called the inner product. Hence, for example, we can define the inner product on the vector space of continuous real-valued function on the interval $[-1,1]$ by 
\begin{equation}
    \bra{f}\ket{g} = \int^{1}_{-1} f(x)g(x)\: dx
\end{equation}
of which then satisfies the inner product axioms. Continuing, we define a few concepts relevant to the space of inner product vector space. 
\begin{definition}[Orthorgonality]
    We say that two vectors $u,v\in E$ are orthogonal if $\bra{u}\ket{v}=0$. 
\end{definition}
\begin{definition}[Norm]
    The normal indicted on a vector space $E$ equipped with an inner product is defined as $||v||=\sqrt{\bra{v}\ket{v}}$ for $v\in E$. 
\end{definition}
\begin{definition}[Orthonormal basis]
    A set of basis vectors all of unit norm ($|u|=1$ for all $u$ of the basis), which are pairwise orthogonal, is called an \textit{orthonormal basis}. 
\end{definition}

The formula for the inner product, component-wise given a coordinate embedded space, is then as followed. Suppose we have two vectors, $\ket{V}$ and $\ket{W}$ expressed by the respective basis, 
\begin{equation}
    \ket{V} = \sum_{i} v_{i}\ket{i},\quad \ket{W} = \sum_{j} w_{j}\ket{j}
\end{equation}
we follow the axioms obeyed by the inner product to obtain 
\begin{equation}
    \bra{V}\ket{W}=\sum_{i}\sum_{j} v_{i}^{*}w_{j}\bra{i}\ket{j}
\end{equation}
where the $\bra{i}\ket{j}$ is the simple dot product between basis vector $i$ and $j$, and the term $v_{i}^{*}w_{j}$ is simply $v_{i}w_{j}$ if we work on real number. Why this equation? Simply because for Euclidean inner product on $\mathbb{F}^{n}$, the inner product formula is 
\begin{equation}
    \bra{(w_{1},w_{2},\dots,w_{n})} \ket{(v_{1},v_{2},\dots,v_{n})} = w_{1}v_{1}^{*} + \dots + w_{n}v^{*}_{n}
\end{equation}
However, there are structures of which this does not hold, hence the proportion is different. There, if $c_{1},\dots,c_{n}$ are positive numbers, of which act like weights, then an inner product can be defined as 
\begin{equation}
    \bra{(w_{1},w_{2},\dots,w_{n})} \ket{(v_{1},v_{2},\dots,v_{n})} = c_{1}w_{1}v_{1}^{*} + \dots + c_{n} w_{n}v^{*}_{n}
\end{equation}
\subsubsection{Gram-Schmidt scheme}

To go any further with the above inner product, we have to know the configuration $\bra{i}\ket{j}$, the inner product between basis vectors. That depends on the details of the basis vectors and all we know for sure is that they are linear independent (which is the criteria of a basis in an inner product space). Suppose that in a two-dimensional space where we have the basis vectors being linearly independent, but not orthogonal. Then, if we write all vectors in terms of this basis, compare to the unit default basis, then we have the inner dot product to be a double sum with four terms, two for each basis vector, as well as the vector components. However, if we use an orthonormal basis such as $\hat{i},\hat{j}$, only diagonal terms like $\bra{i}\ket{i}$ will survive, and we will get only the two terms, without overlapping. Hence that is why we want to have an orthogonal basis, or in general case, orthonormal basis. This is done using \textit{Gram-Schmidt scheme}.
\begin{theorem}[Gram-Schmidt]
    Given a linearly independent basis we can form the linear combinations of the basis vectors to obtain an orthonormal basis. 
\end{theorem}
The detail implementation of such is called the \textit{Gram-Schmidt procedure}, of which we take the virtue to not include here. The inner product then collapse to only $\bra{V}\ket{V}=\sum_{i}v^{*}_{i}w_{i}$, for such basis. 

\subsection{Dual spaces and Dirac notation}

Corresponding to any linear vector space $E$, there exists the dual space of linear functional on $E$. Such space is then called the vector space $E'$ of which for linear functional $F$ assigning a scalar $F(\phi)$ to each vector $phi$, such that to satisfy linearity criteria. To start, we need to know what is a linear functional. A \textit{linear functional}\index{linear functional} on a vector space $E$ is a linear mapping from $E$ to $\mathbb{F}$. From such, we can then define the collection of all such linear functional to a given structure the \textit{dual space}\index{dual space} of all linear functionals on $E$. In other words, $U'=\mathcal{L}(E,\mathbb{F})$. 

Normally, for a linear vector space, we can extrapolate some results and definitions.
\begin{theorem}[Dimensionality]
    Suppose $E$ is finite-dimensional. Then $E'$ is also finite dimensional and $\dim{(E')}=\dim{(E)}$.
\end{theorem} 

The dual space also has its own dual basis, and we will define it in similar fashion as followed (quite similar to a unit basis). 
\begin{definition}[Dual basis]
    If $v_{1},\dots,v_{n}$ is a basis of $E$, then the dual basis of $v_{1},\dots,v_{n}$ is the list $\varphi_{1},\dots,\varphi_{n}$ of elements of $E'$ such that each $\varphi_{j}$ is the linear functional on $E$ such that 
    \begin{equation}
        \varphi_{j}(v_{k}) = \begin{cases}
            1 & k = j \\
            0 & k \neq j
        \end{cases}
    \end{equation}
\end{definition}

The basis works the same for the dual space. Furthermore, we also have that the dual basis of a basis of $E$ consists of the linear functionals on $E$ that give the coefficient for expressing a vector in $E$ as a linear combination of the basis vectors. Quite a connection, but good nonetheless. We refer other concepts, like matrix dual, dual annihilator, dual space properties, dual maps, and so on, to our respective linear algebra chapters.

We note that it is indeed a scalar-inducing space via the inner product on the vector space $E$, if it is equipped with such operator. If $v\in E$, then the map that sends $u$ to $\langle u,v\rangle$ is considered as the linear functional on $E$. How do we know that this is true? Consider the inner product on the inner product space. Such operation is also one kind of linear functional, however, it is encoded into the vector space of such variation itself. What if we can find, for any linear functional in $E'$, an inner product representation of that specific expression? Fortunately, yes, and we leave the credit of such invention to Frigyes Riesz (1880-1956), who proved several theorems just like this, much earlier, and founded this particular representation theorem. Note that by default, linear functional is complete, such is to say they act on for every $u\in E$ of said vector space. 
\begin{theorem}[Riesz representation theorem]
    Suppose $V$ is finite-dimensional and $\varphi$ is a linear functional on $V$. Then, there exist a unique vector $v\in V$ such that 
    \begin{equation}
        \varphi(u) = \langle u,v\rangle
    \end{equation}
    for every $u\in V$. 
\end{theorem}
\begin{proof}
    We first need to show that there exists a vector $v\in V$ such that $\varphi(u)=\langle u,v\rangle$ for every $y\in V$. Let $e_{1},\dots,e_{n}$ be an orthonormal basis of $V$. Then, 
    \begin{align}
        \varphi(u) & = \varphi (\langle \langle u,e_{1} \rangle e_{1} + \dots + \langle  u,e_{n}\rangle e_{n} \rangle)\\
        & =  \langle u,e_{1} \rangle \varphi(e_{1}) + \dots + \langle  u,e_{n}\rangle \varphi(e_{n}) \\
        & = \left\langle u, \varphi^{*}(e_{1})e_{1} + \dots + \varphi^{*}(e_{n}) e_{n}\right\rangle
    \end{align}
    where $\varphi$ has the complex conjugate, also usually denoted as $\bar{\varphi}$. Nevertheless, if it is on real number, then the conjugate is just the real number. Thus setting 
    \begin{equation}
        v= \varphi^{*}(e_{1})e_{1} + \dots \varphi^{*}(e_{n})e_{n}
    \end{equation}
    we have $\varphi(u)=\langle u,v \rangle$ as desired. Now, we need to prove that only vector $v\in V$ has such behaviour. Suppose $v_{1},v_{2}\in V$ are such that 
    \begin{equation}
        \varphi(u) = \langle u,v_{1} \rangle = \langle u,v_{2} \rangle
    \end{equation}
    for every $u\in V$, then 
    \begin{equation}
        0 = \varphi(u) = \langle u,v_{1}\rangle - \langle u,v_{2} \rangle = \langle u,v_{1} - v_{2} \rangle
    \end{equation}
    for every $u\in V$. Taking $u=v_{1}-v_{2}$ shows that $v_{1}-v_{2}=0$. Thus $v_{1}=v_{2}$. This completes the proof of the uniqueness part of the result. 
\end{proof}
\subsection{Operators}
Operator is a step-up from your usual function. Let's say, if you are in Olympiad, probably heard of this in a different context being functional equations. Anyway, moving on, then, an \textit{operator} can be loosely defined as a rule like functions, but instead of variables, you give a function then produce another function outward. Such is to say, given an operator $A$, then $Af(x)=g(x)$ for $g(x)$ is called the \textit{product} of operator $A$ on $f(x)$. Notation varies, but usually it is the uppercase letter, or $\hat{A}$ for example. In physics specifically, of the 3D space, then we have $\hat{A}f(x,y,z,t)=g(x,y,z,t)$ as the functional product. 

As with a lot of things we are concerned of, we are interested in the class of \textit{linear operator} upon an (inner) vector space $V$. Sufficiently saying, denote the object of linear operator as $Q$, then if $\ket{\psi}$ is a vector living in such space, we have 
\begin{equation}
    Q(\alpha \ket{\psi}+ \beta \ket{\chi}) = \alpha(Q\ket{\psi}) + \beta (Q\ket{\chi})
\end{equation}
for $\alpha,\beta\in\mathbb{F}$. In context of quantum mechanics, we change from $\mathbb{F}$ to $\mathbb{C}$. 

To assert equality of two operators, $A=B$, we mean that $A\psi = B\psi$ for all vectors $\psi\in E$. Hence, there exists an isomorphism between them that is satisfied. We can then also define the basic composition rule for them, such as $(A+B)\psi = A\psi + B\psi$, and $AB\psi=A(B\psi)$, and so on. Because they are effectively a functional that acts on function, composition happens, and is written accordingly. For a scalar $k\in \mathbb{C}$, we also have $(kA)f=k(Af)$. Consequently, this means that 
\begin{equation}
    (\alpha \hat{A}+ \beta \hat{B})f(\cdot) = \alpha (\hat{A}f(\cdot)) + \beta 
\end{equation}
The operator can also be related toward the problem of linear algebraical solutions. For an operator $\hat{A}f=g$, there exists a special case that $\hat{A}f=\lambda f$, similar to the problem of eigenvalue and eigenvector. Thereby, instead of eigenvector, the main object of such space is the function space, we would have the concept of \textit{eigenfunction}\index{eigenfunction} $f$ and the eigenvalue $\lambda$ instead. In quantum mechanics, in general, we are interested in this case because with an operator of such kind, it is easier and interpretable of its components to convert from the state to the real space. How can we find the eigenvalue and eigenvector, though? Because we said that operators acts on the function space, then to examine such, we often have to observe how it changes a particular arbitrary function $f$. Let's take an example. 
\begin{example}
    To determine the eigenvalue and eigenfunction of $\hat{A}=i(d/dx)$, we have (suppose there exist such form of the eigenfunction):
    \begin{align}
        \hat{A}f &= \lambda f\\
        i \frac{d}{dx} f &= \lambda f\\
        \frac{d}{dx} f &= -i\lambda f
    \end{align}
    The above equation has the solution as $f=C\exp{(-i\lambda x)}$, and hence we see that $\lambda$ does not depend on $x$. Then there might exist infinitely many eigenvalues.
\end{example}

There are many examples of operators, in general, as their idea is sometimes just generalization of functions. Since they operate for mathematical operations on function, there exists some interesting examples to set. 

\begin{description}[leftmargin = 2.5cm]
    \item[Regular function] A regular function can be thought as  an operator, for $f: x\mapsto ax$, for $a\in \mathbb{R}$. The operator is $f$, operates on real numbers, and multiply the input $x$ by $a$. The function space then reduces to almost as \textit{scalar space}, and hence the function description fits. Squaring is also the same, as $f: x\mapsto x^{2}$. 
    \item[Functional application] For applying on functions, then is the proper terms for operator. $\hat{a}:f(x)\mapsto af(x)$ is an example, where operator $\hat{a}$ operates on scalar function that is mapped by multiplying by $a\in\mathbb{R}$.
    \item[Variable operator] The variable operator is the simple example, for $\hat{x}:f(x)\mapsto xf(x)$. This is the generalization of the above constant multiplication. 
    \item[Differential operator] The operator $D:f(x)\mapsto d/dx f(x)$ is the differential operator, operates on scalar function, which differentiates with respect to $x$ of such function. Partial differentiation operators, gradients, integral operators can be defined similarly. 
    \item[Momentum operator] The momentum operator $\hat{P}$ is defined by $\hat{P}= (\hbar/i)(\partial/\partial x)$, where $\hbar$ is the reduced Planck's constant, $i$ is the imaginary unit, of which operates on the wave function $\psi(x,t)$. For example, it can then be written as \begin{equation}
        \hat{P}(\psi) = \frac{\hbar}{i}\frac{\partial}{\partial x}:\psi(x,t)\mapsto\frac{\hbar}{i}\frac{\partial\psi(x,t)}{\partial x} 
    \end{equation}   
    \item[Energy operator] The energy operator is defined accordingly, as $\hat{E}=(\hbar/i)(\partial/\partial t)$, operates on the wave function. 
    \item[Hamiltonian operator] The Hamiltonian operator is defined, for $m$ is he mass of particle, regarded as constant, to be \begin{equation}
        \hat{H} = - \frac{\hbar^{2}}{2m} \frac{\partial^{2}}{\partial x^{2}}: \psi(x,t)\mapsto  \frac{\hbar^{2}}{2m} \frac{\partial^{2}\psi(x,t)}{\partial x^{2}}
    \end{equation} 
\end{description}

There is also the category of operators which act on vectors, though they are much more complex in notation (still matrix), but usually can be said of such for linear transformation in vector space. 

\subsection{Commutator}

Recalling the concept of an \textit{operator} $A: f^{(n)}(I)\to f(I)$ of which assigns to every function $f\in f^{(n)}(I)$ a function $A(f)\in f(I)$, of the mapping between two function spaces, then for any pair of operators, $A, B$, the \textbf{commutator}\index{commutator} of them is defined as 
\begin{equation}
    \left[ A,B \right] = AB - BA
\end{equation}
The notation here varies, but usually, it is used with the bracket $\left[\cdot\right]$, for operator notation either $A$ or $\hat{A}$ for example (to not mix up stuffs). Simple properties of commutator on operators can be as followed. 

\begin{theorem}[Commutator properties]
    Let $a,b$ be constants, $A,B$ be operators on $n$-differentiable space. Then the following set of equalities hold: 
    \begin{align}
        [f(x),x] & = 0\\
        [A,A] & = 0 \\
        [A,B] & = - [B,A]\\
        [A,BC] & = [A,C]B + A[B,C]\\
        [AB,C] & = [A,C]B + A [B, C]\\
        [a+ A, b+ B] & = [A, B]\\
        [A+B, C+D] &= [A,C] + [A,D] + [B,C] + [B,D]
    \end{align}
\end{theorem}
In general, you can think of the commutator as the interference measurement device. As least to the theory of operator, it transforms the function space, or in general can also be the scalar space, with different ordering of the transformation applied. Remember that operator is indeed a mapping. For case that operator transform on scalar space, for example, $\mathbb{F}^{n}$, then operator reduces to function, and for function space, it maps from this function space to another, basically in both case it works as a transformation on the base space. In general, it basically measures the difference between $AB$, apply $A$ then $B$ of which both acts and transform the space $X\to Y$, then $Y\to Z$, or $BA$, which is the reverse so $X\to Z$ then $Z\to Y$. If they are the same, so basically circular to a result, then the commutator measures zero. This is seen as $[A,A]=0$, as you apply $A$, then continue applying $A$, or by the reverse the path is the same (as connecting two straight lines). We can also then specifically see that it is very much similar to the theory of group, and indeed can be generalized in such direction. The commutator of two group element $A$ and $B$ is $AB^{-1}A^{-1}B^{-1}$, and two elements $A$ and $B$ are said to commute when their commutator is the identity element (same as us). When the group is the \textit{Lie group}, the Lie bracket in its Lie algebra is an infinitesimal version of the group commutator. We defer this to a separate issue on its own, but the idea is the same. 

The last identities to think of aside from those includes the Jacobi's identity, 
\begin{equation}
    [A,[B,C]] + [B,[C,A]] + [C,[A,B]] = 0
\end{equation}
and the generalization of nested operators $[A^n, B]$ etc, 
\begin{equation}
    \left[A,B^{n}\right] = \sum^{n-1}_{j=0}B^{j}[A,B]B^{n-j-1}, \quad \left[A^{n},B\right] = \sum^{n-1}_{j=0} A^{n-j-1}[A,B]A^{j}
\end{equation}
In quantum physics, in general, such expression hence evaluates the interference between two properties of the physical system. If they simply commute, then it is guaranteed for both of them to be measurable on the state space of $\ket{\psi}$ as the wavefunction. 

\subsection{Hermitian operators}

To understand Hermitian operators, we need to understand the adjoints of operators. For now, we are solely in the context of the complex field $\mathbb{C}$, so any kind of adjoint would be between complex operator. By convention, any operator on such function vector space, such as linear operators, can be represented by matrices. Let $H_{1}$ and $H_{1}$ be Hilbert space $X$ over a field $\mathbb{F}\in\{\mathbb{R},\mathbb{C}\}$. We define the adjoint\index{adjoint operator} as followed:
\begin{theorem}[Adjoint operator]
    Let $T\in\mathcal{B}(X)$ be a bounded linear operator on a Hilbert space $X$. There exists a unique operator $T^{*}\in\mathcal{B}(X)$ such that 
    \begin{equation}
        \langle T x, y\rangle = \langle x, T^{*}y\rangle, \quad \forall x,y\in X
    \end{equation}
    The operator $T^{*}$ is called the \textbf{adjoint} of $T$. 
\end{theorem}
This is true for also finite dimensional inner product space, and so is Hilbert space through assumptions on Riesz representation theorem. The notion of adjoint is the functional generalization to the concept of transpose. By the definition of transpose, a transpose on a matrix $X$ as a linear transformation reverses the direction of said transformation on the inner product space per definition of said space. In a sense, you can think of it as transforming perspective, from linear operations on $\mathbb{V}$ into the space of covectors, or dual space $\mathbb{V}'$. Remember that the dual space represents measurements and functional application on vector space. Hence, basically, transpose do the operations on $W^{*}\to V^{*}$ for a linear transformation $V\to W$. This works on the complex vector space, where the inner product involves complex conjugation, and in such case there exists the \textbf{adjoint matrix} for a given complex transformation. 

Something is then called a \textbf{Hermitian operator} if the adjoint is the operator, or $T^{*}=T$. Quantum mechanically, for $\psi, \phi$ being arbitrary functions on complex space. Then, if operator $\hat{A}$ is such that \index{Hermitian operator}
\begin{equation}
    \int \psi (\hat{A}\phi)^{*}\: d\tau = \in t \psi^{*}(\hat{A}\psi)\: d\tau
\end{equation}
then it is called a quantum Hermitian operator, or \textit{self-adjoint} operator. Under matrix notation, then adjoint is the \textit{conjugate transpose} of such complex matrix. We have the following definition for \textit{skew-adjoint} and \textit{self-adjoint} operators:
\begin{definition}[Skew-adjoint operators]\index{skew-adjoint}
    A linear operator $A: \mathcal{U}\to\mathcal{Q}$ is said to be \textit{skew-adjoint} if $A^{*}=-A$, that is 
    \begin{equation}
        \langle A\mathbf{u}, \mathbf{v}\rangle = -\langle \mathbf{u} , A\mathbf{V}\rangle
    \end{equation}
    for all $\mathbf{u,v}\in \mathcal{U}$. 
\end{definition}
\begin{definition}[Self-adjoint operators]\index{self-adjoint}
    A linear operator $A: \mathcal{U}\to\mathcal{Q}$ is said to be \textit{self-adjoint} if $A^{*}=A$, that is 
    \begin{equation}
        \langle A\mathbf{u}, \mathbf{v}\rangle = \langle \mathbf{u} , A\mathbf{V}\rangle
    \end{equation}
    for all $\mathbf{u,v}\in \mathcal{U}$. 
\end{definition}
All of these definitions are applicable for both real and complex space, of the metric standard. 
\subsubsection{A closer look at complex adjoint}

Previously, we kind of mix up between the field $\mathbb{F}$, where it is arbitrary of the complex or real space. For now, let us examine such adjoint in the complex space itself. 

Given the linear mapping $A:\mathbb{C}^{n}\to \mathbb{C}^{m}$, then it is easiest to think of 
\begin{equation}
    A\mathbf{u} = u_{1} \mathbf{a}_{1} + \dots u_{n}\mathbf{a}_{n}
\end{equation}
for the transformation basis $\mathbf{a}_{i}$. Thus, we may deduce the $i$th entry of $A\mathbf{u}$ given by 
\begin{equation}
    (A\mathbf{u})_{i} = u_{1}(\mathbf{a}_{1})_{i}+ \dots u_{n}(\mathbf{a}_{n})_{i} = \sum^{n}_{j=1}u_{j}a_{ij}
\end{equation}
We define the complex version of the inner product as the \textbf{Hermitian inner product}, such as followed. 
\begin{definition}[Complex inner product]
    The vector space $V$ over field $\mathbb{C}$ is said to have the complex inner product $\langle\cdot,\cdot\rangle:V\times V\to \mathbb{C}$ if and only if it satisfies, for all $x,y,y_{1},y_{2}\in V$ and $\alpha\in \mathbb{C}$ the following set of conditions:
    \begin{align}
        \langle x, y_{1}+y_{2}\rangle &= \langle x,y_{1}\rangle + \langle x, y_{2}\rangle \\
        \langle x, \alpha y\rangle & = \alpha \langle x,y \rangle\\
        \langle y, x \rangle &= \langle x,y\rangle^{*}\\
        \langle x, x \rangle &> 0 \: \text{if}\: x\ne 0
    \end{align}
\end{definition}

Thus, if $(A\mathbf{u})_{i} =\sum^{n}_{j=1}u_{j}a_{ij}$, then we have the inner product as 
\begin{equation}
    \begin{split}
        \langle A\mathbf{u},\mathbf{v} \rangle 
        & = \sum^{m}_{i=1} \left(\sum^{n}_{j=1}(u_{j}a_{ij})\right)^{*}v_{i}\\
        & = \sum^{m}_{i=1} \left(\sum^{n}_{j=1}(u_{j}^{*}a_{ij}^{*})v_{i}\right)\\
        & = \sum^{n}_{j=1}\left(\sum^{m}_{i=1}(u_{j}^{*}a_{ij}^{*})v_{i}\right)\\
        & = \sum^{n}_{j=1}\left(\sum^{m}_{i=1}(v_{i}a_{ij}^{*})u^{*}_{j}\right)\\
        & = \sum^{n}_{j=1}\left(\sum^{m}_{i=1}(v_{i}a_{ij}^{*})\right)u^{*}_{j} \\
        & = \langle \mathbf{u}, A^{*}\mathbf{v}\rangle
    \end{split}
\end{equation}

Here, the $ji$ entry of $A^{*}$ is the conjugate of the $ij$ entry of $A$. How do we interpret such result? Basically, here we gain \begin{equation}
    (A^{*}\mathbf{v})_{i} = \sum^{m}_{i=1} (v_{i}a_{ij}^{*}), \quad j = 1,\dots, n
\end{equation}
Of which looks like the dual form of $(A\mathbf{u})_{i}$. And indeed it is so much of the change in basis, especially when we have the decomposition: 
\begin{equation}
    (A^{*}\mathbf{v}) = \sum^{m}_{i} v_{i}a^{*}_{ij} = v_{1} a^{*}_{1j} + \dots v_{m}a_{mj}^{*} 
\end{equation}
This is the linear combination of the conjugates of the rows of $A$. In general, if a theorem refers to the adjoint of a linear map, then if you are applying that to linear maps on finite dimensional complex vector spaces, you can equivalently substitute that with the phrase conjugate transpose. The conjugate transpose interact with the inner product as the general adjoint is. For a linear map $A:\mathbb{C}^{n}\to\mathbb{C}^{m}$, then $A$ can be represented as an $m\times n$ complex matrix. The adjoint by then is the $n\times m$ matrix that is the matrix $A$ reflected through the diagonal with each entry conjugated, and we call it conjugate transpose. 
\subsubsection{Quantum Hermitian operators}
As we have said, the operator $P$ is defined as Hermitian if it satisfies 
\begin{equation}
    \int \psi^{*}(x) P \phi(x) \: dx = \int (P \psi(x))^{*}\phi(x)\: dx
\end{equation}
To see how this definition works, let's show that the momentum operator, $p=-\hbar d/dx$ is Hermitian. We have: 
\begin{align}
    \int \psi^{*}(x) \left(-i\hbar \frac{d}{dx}\right)\phi(x)\: dx 
    & = -i\hbar \int \psi^{*}(x)\frac{d}{dx}\phi(x)\: dx \\
    & = -i\hbar \int \psi(x)^{*} d\phi(x)\\
    & = -i\hbar \left\{ \psi^{*}(x)\phi(x) - \int \phi(x)d\psi^{*}(x) \right\}\\
    & = \int \phi(x)\left(-i\hbar \frac{d}{dx}\psi(x)\right)^{*}
\end{align}
which is a Hermitian by definition of hermiticity. There are three important consequences of an operator being hermitian. Its eigenvalues are \textit{real}, its eigenfunction corresponding to different eigenvalues are orthogonal to each other, and the set of all its eigenfunctions is complete. Observations in reality is always measured of real number, hence we can interpret the eigenvalues as observable quantities of a given quantum system (we will get to this seriously later on). Eigenfunctions being orthogonal basically allows us, to get exclusive state. Hence, across all states, you can write a general state $\Phi$ as the combination of the orthogonal states as 
\begin{equation}
    \Phi = c_{1}\phi_{1} + c_{2}\phi_{2} + \dots 
\end{equation}
for the $\{\phi_{i}\}$ combinations. Completeness basically means that eigenfunctions here form a complete basis of the state space, such that any $\Phi$ can be represented with just as many combinations. 